\section{Conclusion}

A balanced update is straightforward. Bard et al.\ were right to insist that
acceptability could be treated as measurable and methodologically serious. Their documented case
for magnitude estimation concerned scale resolution, fine distinctions,
naive-participant use, and replication across groups
\citep[32]{BardEtAl1996}. That claim deserves credit as part of the
history of experimental syntax.

Method-specific conclusions require more restraint and a fair test.
Magnitude estimation is compared against a measurement null in which later
formal tasks read a common item-level acceptability signal for the analysed
contrasts through method-specific response functions, and in which the
single-dimension assumption is checked before any method-specific term is
read. The question is then whether magnitude estimation earns its keep on
resolution, the distinctions Bard argued fixed scales censor, or on prediction,
decision stability, or construct interpretation, rather than whether it survives
a null that books resolution as noise.

The approximation uses aggregate matrices from the two comparison studies,
response-function diagnostics, pair-level robustness checks, and a bounded post-outcome
multiverse. Its scope matters because the participant rows needed for a full
dimensionality model across Bard, Sprouse, and Langsford aren't available to
this project.

The reanalysis supports that update. Across the two comparison datasets,
magnitude estimation and the bounded methods share a strong aggregate practical
signal. First-pass aggregate
dimensionality checks show no obvious evidence against a single dominant
dimension at that level of aggregation. Raw Likert means show bounded
response-function curvature, which is exactly why the null shouldn't treat
bounded tasks as simple noisy copies of magnitude estimation.

By contrast, the remaining resolution diagnostics are less favourable to the
stronger method claim: endpoint regions don't show unusually large magnitude-estimation
spread, and pair-level bounded predictors recover most magnitude-estimation
contrast variance. None of the three multiverse families clears its frozen
cross-specification criterion in either dataset. Four 2017 rank-scale
forced-choice specifications show a small ME increment, but that result doesn't
generalize to yes/no judgments, raw-scale mappings, or the 2013 data. Langsford
et al.\ remain article-level evidence because their raw responses aren't
available to this project.

Projectibility adds a correspondingly modest payoff. Within the sampled
contrasts in the two comparison datasets, observing relative ordering under one
formal task warrants expecting similar aggregate ordering under another.
Magnitude estimation alone doesn't
support an additional robust practical ordering for those contrasts. The result
warrants a task-format inference, not the claim that every grammatical inference travels
across speakers, registers, or observation channels.

The conclusion would change under a clear ME-only pattern: endpoint regions
with much larger ME spread than middle regions, high-ME contrasts that bounded
methods consistently placed low, stable method-specific residual structure
across scoring choices, or lower sign-error and exaggeration risk for named
linguistic contrasts. Those are the patterns the present diagnostics do
not show.

Across these checks, the central conclusion remains stable. The durable legacy is the formal
measurement of graded acceptability as one route into grammar's projectibility
profiles. The broader privileged status sometimes read into magnitude
estimation is a separate claim, and later evidence has to earn it separately.
